% GENERATED FILE. Edit canonical lessons, then run npm run generate:books.
% canonical-source-hash: fnv1a64:1b5029b8ebbc61df
% canonical-lessons: PT-C03-eu, PT-C03-me-chamo, PT-C03-como-se-chama, PT-C03-prazer, PT-C03-practice

\chapter{Introducing Yourself}
\label{ch:pt-introductions}
\hlchaptermodality{3}

\begin{chapteropening}
\textbf{By the end of this chapter:} \emph{I can give my name in Portuguese, ask for someone else's, and say that it is a pleasure.}

A whole first meeting: Olá, me chamo …, Como se chama?, prazer — then the same exchange again with o senhor and a senhora.
\end{chapteropening}

\section[eu]{eu — “I,” often dropped in Portuguese}

\label{lesson:PT-C03-eu}



\begin{quote}
Before introducing yourself, meet \textbf{eu} — \textquotedblleft{}I.\textquotedblright{} Famous English cousin, and the same habit as its Romance sisters: Portuguese usually \textbf{leaves it out.}
\end{quote}

\begin{sounds}
\begin{itemize}
\raggedright
  \item \texttt{nasal-eu} — \textbf{eu} = \emph{EH-oo}, a gliding vowel with a light nasal colour, one syllable.
\end{itemize}
\end{sounds}

\begin{cousinweb}
\textbf{eu} = \textquotedblleft{}I,\textquotedblright{} from Latin \textbf{ego} — the same \emph{ego} English borrowed outright:

\begin{itemize}
\raggedright
  \item \textbf{ego}, \textbf{egoist}, \textbf{egocentric} — all Latin \textquotedblleft{}I.\textquotedblright{}
  \item Romance siblings, all from \emph{ego}: Spanish \textbf{yo}, Italian \textbf{io}, French \textbf{je}.
\end{itemize}
\end{cousinweb}

\begin{grammarlens}[title={Grammar Lens: the dropped subject (pro-drop)}]
The verb ending already names the speaker, so Portuguese usually \textbf{omits} the pronoun: \emph{chamo-me} (\textquotedblleft{}I call myself\textquotedblright{}) needs no \emph{eu} — the \emph{-o} of \emph{chamo} is \textquotedblleft{}I.\textquotedblright{} You add \textbf{eu} for \textbf{emphasis} or \textbf{contrast}:

\begin{quote}
\emph{Eu} estou bem, e você? — \textquotedblleft{}\emph{I'm} fine, and you?\textquotedblright{}
\end{quote}

Spanish (\emph{yo}) and Italian (\emph{io}) do the same. It's one of the first real habits of a Romance language: trust the verb ending, drop the pronoun.
\end{grammarlens}

\subsection*{Your turn}
Say these aloud:

\begin{itemize}
\raggedright
  \item \textquotedblleft{}eu\textquotedblright{} — \emph{EH-oo}
  \item \textquotedblleft{}eu\textquotedblright{} then English \textquotedblleft{}ego, egoist\textquotedblright{} — the same Latin \textquotedblleft{}I\textquotedblright{}
  \item eu / yo / io / je — four children of \emph{ego}
\end{itemize}

\begin{tcolorbox}[breakable,colback=teal!4,colframe=teal!35!black,title={Before you move on}]
What Latin word is \emph{eu} from? (\emph{Ego} — ego, egoist.) Why does Portuguese usually drop it? (The verb ending already says \textquotedblleft{}I.\textquotedblright{}) When do you keep it? (Emphasis / contrast.) Next: introduce yourself — \textbf{me chamo}.
\end{tcolorbox}
\section[me chamo / chamo-me]{me chamo — \textquotedblleft{}I call myself\textquotedblright{}}

\label{lesson:PT-C03-me-chamo}



\begin{quote}
Like its Iberian and Italian sisters, Portuguese introduces you as what you \textbf{call yourself} — \textbf{me chamo} — from the very same Latin \emph{clāmāre} that gave Spanish \emph{me llamo} and Italian \emph{mi chiamo}.
\end{quote}

\begin{sounds}
\begin{itemize}
\raggedright
  \item \texttt{ch-is-sh} — Portuguese \textbf{ch} = \textbf{\textquotedblleft{}sh\textquotedblright{}}: \textbf{chamo} = \emph{SHAH-moo}. (A neat three-way split from one Latin \emph{cl-}: Spanish \emph{ll} = \emph{y} (\emph{llamo}), Italian \emph{ch} = hard \emph{k} (\emph{chiamo}), Portuguese \emph{ch} = \emph{sh} (\emph{chamo}).)
  \item \texttt{final-o-is-u} — final \emph{-o} $\to$ \emph{u}: \emph{SHAH-moo}.
\end{itemize}
\end{sounds}

\begin{cousinweb}
\textbf{me chamo} = \textbf{me} (\textquotedblleft{}myself\textquotedblright{}) + \textbf{chamo} (\textquotedblleft{}I call\textquotedblright{}) $\to$ \textbf{\textquotedblleft{}I call myself.\textquotedblright{}} Brazilians often say \textbf{eu me chamo}; European Portuguese prefers \textbf{chamo-me} (same words, pronoun clipped to the back) — both fine. Then the name: \emph{Me chamo Ana} — \textquotedblleft{}I call myself Ana.\textquotedblright{}

The verb \textbf{chamar-se} (\textquotedblleft{}to call oneself\textquotedblright{}) is from Latin \textbf{clāmāre}, \emph{\textquotedblleft{}to cry out, to call\textquotedblright{}} — shouting through English as \textbf{claim}, \textbf{exclaim}, \textbf{proclaim}, \textbf{clamor}. Three Romance naming verbs, one Latin call:

\noindent
\begin{tabularx}{\linewidth}{@{}>{\raggedright\arraybackslash}X>{\raggedright\arraybackslash}X>{\raggedright\arraybackslash}X@{}}
\toprule
\textbf{language} & \textbf{\textquotedblleft{}I call myself\textquotedblright{}} & \textbf{\emph{cl-} became} \\
\midrule
\textbf{Portuguese} & \emph{me chamo} & \emph{ch} (= \emph{sh}) \\
\textbf{Spanish} & \emph{me llamo} & \emph{ll} (= \emph{y}) \\
\textbf{Italian} & \emph{mi chiamo} & \emph{ch} (= hard \emph{k}) \\
\bottomrule
\end{tabularx}

\textbf{Another way to say it:} \emph{o meu nome é …} (\textquotedblleft{}my name is …\textquotedblright{}), where \textbf{nome} $\leftarrow$ Latin \textbf{nōmen}, \textquotedblleft{}name\textquotedblright{} — the root of English \textbf{noun}, \textbf{nominal}, \textbf{nomenclature}, \textbf{nominate}. Two routes to the same introduction: the \emph{verb} (call myself) or the \emph{noun} (my name is).
\end{cousinweb}

\begin{grammarlens}[title={Grammar Lens: the reflexive}]
\emph{chamo} alone = \textquotedblleft{}I call {[}someone{]}\textquotedblright{}; \textbf{me} chamo = \textquotedblleft{}I call \textbf{myself}\textquotedblright{} — reflexive. The pronoun shifts with the person: \textbf{me} chamo (I) $\to$ \textbf{se} chama (he/she/você). And \emph{eu} is dropped — the \emph{-o} ending already says \textquotedblleft{}I.\textquotedblright{}
\end{grammarlens}

\subsection*{Your turn}
Say these aloud:

\begin{itemize}
\raggedright
  \item \textquotedblleft{}me chamo\textquotedblright{} — \emph{muh SHAH-moo}, \emph{ch} = \emph{sh}
  \item \textquotedblleft{}Me chamo …\textquotedblright{} with your own name
  \item me chamo / me llamo / mi chiamo — one \emph{clāmāre}, three sounds
\end{itemize}

\begin{tcolorbox}[breakable,colback=teal!4,colframe=teal!35!black,title={Before you move on}]
What does \emph{me chamo} literally mean, and its Latin root? (\textquotedblleft{}I call myself\textquotedblright{}; \emph{clāmāre} — claim/clamor.) How is Portuguese \emph{ch} pronounced, vs Italian \emph{ch}? (Portuguese \emph{sh}; Italian hard \emph{k}.) What's the noun-route alternative? (\emph{O meu nome é …}, \emph{nome} $\leftarrow$ \emph{nōmen}.) Next: ask it back — \textbf{Como se chama?}
\end{tcolorbox}
\section[Como se chama?]{Como se chama? — \textquotedblleft{}what's your name?\textquotedblright{}}

\label{lesson:PT-C03-como-se-chama}



\begin{quote}
Both pieces are yours: \textbf{como} (\textquotedblleft{}how,\textquotedblright{} from your \emph{Tudo bem?} chapter) and \textbf{chamar-se} (\textquotedblleft{}to call oneself\textquotedblright{}). Portuguese, like all its sisters, asks the name with \textbf{\textquotedblleft{}how\textquotedblright{}}: \emph{how do you call yourself?}
\end{quote}

\begin{cousinweb}
\begin{quote}
\textbf{Como se chama?} = \textquotedblleft{}What's your name?\textquotedblright{} literally: \textbf{\textquotedblleft{}How do you call yourself?\textquotedblright{}}
\end{quote}

\begin{itemize}
\raggedright
  \item \textbf{como} = \textquotedblleft{}how\textquotedblright{} ($\leftarrow$ \emph{quōmodo}).
  \item \textbf{se} = \textquotedblleft{}yourself\textquotedblright{} (the reflexive; \emph{me} $\to$ \emph{se}).
  \item \textbf{chama} = \textquotedblleft{}calls\textquotedblright{} (\emph{você}/he-she form of \emph{chamar-se}).
\end{itemize}

Portuguese register is friendly here — the everyday polite form uses \textbf{você}:

\noindent
\begin{tabularx}{\linewidth}{@{}>{\raggedright\arraybackslash}X>{\raggedright\arraybackslash}X@{}}
\toprule
\textbf{form} & \textbf{question} \\
\midrule
polite / neutral (\emph{você}) & \textbf{Como (você) se chama?} \\
very formal & \textbf{Como se chama o senhor / a senhora?} \\
familiar (\emph{tu}, more in Portugal) & \textbf{Como te chamas?} \\
\bottomrule
\end{tabularx}

The reply loops right back: \emph{Me chamo Ana. E você?} And just like Spanish \emph{¿cómo se llama?} and Italian \emph{come ti chiami?}, the name is a matter of \textbf{how} you're called — never \emph{what}.
\end{cousinweb}

\subsection*{Your turn}
Say these aloud:

\begin{itemize}
\raggedright
  \item \textquotedblleft{}Como se chama?\textquotedblright{} — \emph{KOH-moo suh SHAH-mah}
  \item the whole loop — \textquotedblleft{}Como se chama? — Me chamo … E você?\textquotedblright{}
  \item como se chama / cómo se llama / come ti chiami — the same question, three sisters
\end{itemize}

\begin{tcolorbox}[breakable,colback=teal!4,colframe=teal!35!black,title={Before you move on}]
What does \emph{Como se chama?} literally ask? (\textquotedblleft{}How do you call yourself?\textquotedblright{}) Which everyday pronoun does the polite form use? (\emph{Você}.) What \textquotedblleft{}how\textquotedblright{} word is it built on? (\emph{Como} $\leftarrow$ \emph{quōmodo}.) Next: the warm reply — \textbf{prazer}.
\end{tcolorbox}
\section[prazer]{prazer — \textquotedblleft{}pleased to meet you,\textquotedblright{} i.e. \textquotedblleft{}a pleasure\textquotedblright{}}

\label{lesson:PT-C03-prazer}



\begin{quote}
Names traded, you close with \textbf{prazer} — \textquotedblleft{}a pleasure.\textquotedblright{} It's the twin of Italian \emph{piacere}, and, like it, straight from Latin \emph{please}.
\end{quote}

\begin{sounds}
\begin{itemize}
\raggedright
  \item \texttt{pr-cluster} — \textbf{prazer} = \emph{prah-\textbf{ZEHR}}: the \emph{z} is a true \emph{z}; stress the end. (The \emph{pr-} shows a Portuguese sound-law: Latin \emph{pl-} often became \emph{pr-} — \emph{placēre} $\to$ \emph{prazer}, \emph{plūs} $\to$ …)
\end{itemize}
\end{sounds}

\begin{cousinweb}
\textbf{prazer} = \textquotedblleft{}pleasure\textquotedblright{} (and, as a verb, \textquotedblleft{}to please\textquotedblright{}), from Latin \textbf{placēre}, \emph{\textquotedblleft{}to please.\textquotedblright{}} On meeting someone it's short for \emph{\textquotedblleft{}{[}it's{]} a pleasure.\textquotedblright{}}

That \emph{placēre} root is deep in English — \textbf{please}, \textbf{pleasure}, \textbf{pleasant}, \textbf{complacent}, \textbf{placid} — and it names this courtesy across the family:

\begin{itemize}
\raggedright
  \item Italian \textbf{piacere} — its exact twin (same \emph{placēre}).
  \item Spanish \textbf{mucho gusto} — \textquotedblleft{}much pleasure,\textquotedblright{} but from \emph{gustus} (\textquotedblleft{}taste\textquotedblright{}), a different root.
  \item French \textbf{enchanté} — a different picture again (\textquotedblleft{}enchanted\textquotedblright{}).
\end{itemize}

Portuguese took \emph{placēre}'s \emph{pl-} to \emph{pr-} (a regular shift — compare \emph{branco} \textquotedblleft{}white\textquotedblright{} $\leftarrow$ Germanic \emph{blank}, \emph{praia} \textquotedblleft{}beach\textquotedblright{} $\leftarrow$ \emph{plagia}), so \emph{piacere} and \emph{prazer} look a little different but are the very same word.
\end{cousinweb}

\begin{grammarlens}[title={Grammar Lens: a one-word courtesy}]
\emph{Prazer} stands alone, with a nod or handshake. Fuller forms: \textbf{muito prazer} (\textquotedblleft{}much pleasure\textquotedblright{}), or \textbf{prazer em conhecê-lo/-la} (\textquotedblleft{}a pleasure to know you\textquotedblright{}).
\end{grammarlens}

\subsection*{Your turn}
Say these aloud:

\begin{itemize}
\raggedright
  \item \textquotedblleft{}prazer\textquotedblright{} — \emph{prah-ZEHR}
  \item \textquotedblleft{}prazer\textquotedblright{} then English \textquotedblleft{}please, pleasure, placid\textquotedblright{} — the \emph{placēre} root
  \item Portuguese \emph{prazer} and Italian \emph{piacere} — twins from \emph{placēre}
\end{itemize}

\begin{tcolorbox}[breakable,colback=teal!4,colframe=teal!35!black,title={Before you move on}]
What does \emph{prazer} mean, and from what Latin verb? (\textquotedblleft{}A pleasure,\textquotedblright{} $\leftarrow$ \emph{placēre}.) Its Italian twin? (\emph{Piacere}.) What sound-shift turned \emph{pl-} into \emph{pr-}? (The Portuguese \emph{placēre} $\to$ \emph{prazer} change.) Next: put the whole introduction together.
\end{tcolorbox}
\section[Practice]{Practice — introducing yourself in Portuguese}

\label{lesson:PT-C03-practice}



\begin{quote}
No new words. \emph{eu, me chamo, como se chama, prazer} assemble into a real first meeting. This slots between your greetings (Ch. 1–2) and farewells (Ch. 4), so Portuguese now runs a whole conversation.
\end{quote}

\subsection*{The exchange}
\textbf{Casual} (\emph{você}):

\begin{quote}
— Olá! \textbf{Como você se chama?} — \textbf{Me chamo} Ana. E você? — João. \textbf{Prazer!} — Prazer!
\end{quote}

\textbf{More formal} (o senhor / a senhora):

\begin{quote}
— Bom dia. \textbf{Como se chama a senhora?} — Chamo-me Ana Costa. E o senhor? — Costa. \textbf{Muito prazer.}
\end{quote}

\begin{grammarlens}[title={What you've built this chapter}]
\begin{itemize}
\raggedright
  \item \textbf{eu} — \textquotedblleft{}I\textquotedblright{} ($\leftarrow$ \emph{ego}), usually \textbf{dropped} — the verb ending carries it.
  \item \textbf{me chamo / chamo-me} — \textquotedblleft{}I call myself\textquotedblright{} (\emph{chamar-se} $\leftarrow$ \emph{clāmāre}; \emph{ch} = \emph{sh}; kin of \emph{me llamo}, \emph{mi chiamo}). Alt: \emph{o meu nome é} (\emph{nome} $\leftarrow$ \emph{nōmen}).
  \item \textbf{Como se chama?} — \textquotedblleft{}how do you call yourself?\textquotedblright{} (polite \emph{você}).
  \item \textbf{prazer} — \textquotedblleft{}a pleasure\textquotedblright{} ($\leftarrow$ \emph{placēre}; twin of Italian \emph{piacere}).
\end{itemize}
\end{grammarlens}

\subsection*{Your turn}
Say these aloud:

\begin{itemize}
\raggedright
  \item the full casual meeting, both voices
  \item the same formally — with \emph{o senhor / a senhora}
  \item chain it all — greet, name, ask, how-are-you, goodbye: \textquotedblleft{}Olá! Me chamo … Como se chama? … Tudo bem? … Até logo!\textquotedblright{}
\end{itemize}

\emph{Twice through:} Run a whole first meeting start to finish.

\begin{tcolorbox}[breakable,colback=teal!4,colframe=teal!35!black,title={Before you move on}]
Why does Portuguese usually drop \emph{eu}? (The verb ending says who.) What does \emph{me chamo} literally mean, and its root? (\textquotedblleft{}I call myself\textquotedblright{}; \emph{clāmāre}.) What does \emph{prazer} mean, and its Italian twin? (\textquotedblleft{}A pleasure\textquotedblright{}; \emph{piacere}.) Portuguese now runs greet $\to$ introduce $\to$ how-are-you $\to$ goodbye, end to end — the loop is closed.
\end{tcolorbox}
